\documentclass[../principal.tex]{subfiles}
\graphicspath{{\subfix{../imagenes/}}}

\begin{document}

\chapter{Condicionales}

\section{Definición}

Las condicionales condicionan.

Las condicionales permiten que nuestro código se bifurque, que tome distintos caminos dependiendo de condiciones.

Las condicionales también son conocidas en inglés como conditionals, o if statements.

\section{Sintaxis}

Las condicionales en JavaScript se declaran usando la palabra reservada $if$.

Luego de la palabra reservada $if$ se escribe una pregunta entre paréntesis, y luego un bloque de código entre llaves.

\begin{lstlisting}[language=JavaScript]
if (pregunta) {
    // código a correr si la respuesta es afirmativa
}
\end{lstlisting}

\end{document}


